\GalSourcePageBoundary{068}{L0067}{39}{39}

les racines soient fonctions rationnelles de $V$. Cela posé, considérons
l'équation irréductible dont $V$ est racine (lemmes III et IV).
Soient $V,V',V'',\ldots,V^{(n-1)}$ les racines de cette équation.

Soient $\varphi V,\varphi_{1}V,\varphi_{2}V,\ldots,\varphi_{m-1}V$ les racines de la proposée.

Écrivons les $n$ permutations suivantes des racines:
\[
\begin{array}{c@{\quad}ccccc}
(V)&\varphi V&\varphi_{1}V&\varphi_{2}V&\ldots&\varphi_{m-1}V,\\
(V')&\varphi V'&\varphi_{1}V'&\varphi_{2}V'&\ldots&\varphi_{m-1}V',\\
(V'')&\varphi V''&\varphi_{1}V''&\varphi_{2}V''&\ldots&\varphi_{m-1}V'',\\
\ldots&\ldots&\ldots&\ldots&\ldots&\ldots,\\
(V^{(n-1)})&\varphi V^{(n-1)}&\varphi_{1}V^{(n-1)}&\varphi_{2}V^{(n-1)}&\ldots&\varphi_{m-1}V^{(n-1)}:
\end{array}
\]
je dis que ce groupe de permutations jouit de la propriété énoncée.

En effet:

1\textsuperscript{o} Toute fonction $F$ des racines, invariable par les substitutions
de ce groupe, pourra être écrite ainsi: $F=\psi V$, et l'on aura
\[
\psi V=\psi V'=\psi V''=\ldots=\psi V^{(n-1)}.
\]
La valeur de $F$ pourra donc se déterminer rationnellement.

2\textsuperscript{o} Réciproquement, si une fonction $F$ est déterminable rationnellement,
et que l'on pose $F=\psi V$, on devra avoir
\[
\psi V=\psi V'=\psi V''=\ldots=\psi V^{(n-1)},
\]
puisque l'équation en $V$ n'a pas de diviseur commensurable et que
$V$ satisfait à l'équation $F=\psi V$, $F$ étant une quantité rationnelle.
Donc la fonction $F$ sera nécessairement invariable par les substitutions
du groupe écrit ci-dessus.

Ainsi, ce groupe jouit de la double propriété dont il s'agit dans
le théorème proposé. Le théorème est donc démontré.

Nous appellerons \emph{groupe de l'équation} le groupe en question.

\emph{Scolie I.} --- Il est évident que, dans le groupe de permutations
dont il s'agit ici, la disposition des lettres n'est point à considérer,
mais seulement les substitutions de lettres par lesquelles on passe
d'une permutation à l'autre.

Ainsi l'on peut se donner arbitrairement une première permutation,
pourvu que les autres permutations s'en déduisent toujours
par les mêmes substitutions de lettres. Le nouveau groupe ainsi
formé jouira évidemment des mêmes propriétés que le premier,
